\documentclass[11pt,a4paper]{article}

\usepackage[utf8]{inputenc}
\usepackage[T1]{fontenc}
\usepackage{amsmath,amssymb,mathtools}
\usepackage{geometry}
\usepackage{hyperref}
\usepackage{booktabs}
\usepackage{array}
\usepackage{enumitem}

\geometry{margin=1in}
\numberwithin{equation}{section}
\setlist[itemize]{itemsep=0.2em,topsep=0.3em}
\setlist[enumerate]{itemsep=0.2em,topsep=0.3em}

\title{\textbf{Yee-Cell Finite-Difference Solver for 2D Electromagnetics}\\[0.4em]
\large Formulation, Discretization, and Implementation Plan\\
with a Roadmap to GSTC Metasurfaces}
\author{Tom Smy}
\date{\today}

\begin{document}
\maketitle

\begin{abstract}
This note is the specification for MaxwellFD, a Python finite-difference code
on a uniform Cartesian Yee grid. One grid, one material sampling, and one pair
of discrete curls serve both the leapfrog scheme (FDTD) and the time-harmonic
sparse scheme (FDFD). Phasors use $e^{j\omega t}$, the same convention as the
sibling surface code MaxwellMOM. The derivation below is two-dimensional, for
TMz and TEz, in free space and in a uniform isotropic medium. PEC and periodic
walls, the CFL limit, numerical dispersion, a stretched-coordinate PML, and a
resistive sheet are written out against equation numbers. The resistive
conductance is in the code. The non-dispersive GSTC subcell is the symmetric
Yee cell: the face unknowns and the local system are derived here, and the
tangential sheet is in the code. Spatially varying $\chi$, Bloch cells, and
dispersive $\chi$ are specified as checks, not as this slice's code.
\end{abstract}

\tableofcontents
\newpage

\section{Mathematical formulation}
\label{sec:formulation}

\subsection{Maxwell's equations}
\label{sec:maxwell}

Fields are in SI units. The time convention is $e^{j\omega t}$, so a wave
$e^{-jk\hat{\mathbf{k}}\cdot\mathbf{r}}$ propagates in the direction
$\hat{\mathbf{k}}$ and $\partial_t$ becomes $j\omega$. With an impressed
current $\mathbf{J}$,
\begin{align}
\nabla\times\mathbf{E} &= -j\omega\mu\,\mathbf{H}, \label{eq:faraday}\\
\nabla\times\mathbf{H} &= j\omega\varepsilon\,\mathbf{E}+\mathbf{J}. \label{eq:ampere}
\end{align}
Conduction is not a second current on the right. It stays in a real
conductivity $\sigma$ that the time stepper absorbs into its coefficients, or,
in a later frequency-domain material, in a complex permittivity. The
constants $c_0$, $\varepsilon_0$, $\mu_0$, and $\eta_0$ are the SI values used
by MaxwellMOM. In a uniform lossless medium $c=1/\sqrt{\mu\varepsilon}$ and
$\eta=\sqrt{\mu/\varepsilon}$.

The impressed current is the source. A plane wave, a dipole, and a Gaussian
beam are later special cases of $\mathbf{J}$ or of a total-field/scattered-field
boundary. Milestone~1 does not inject those sources. Its free-space check is a
periodic eigenmode, and its cavity check is a source-free resonance.

\subsection{Two-dimensional reductions}
\label{sec:polarizations}

The grid lies in the $xy$-plane and nothing depends on $z$.

\paragraph{TMz.}
The unknowns are $E_z$, $H_x$, and $H_y$. Equations~\eqref{eq:faraday}
and~\eqref{eq:ampere} reduce to
\begin{align}
\partial_y E_z + j\omega\mu H_x &= 0, \label{eq:tm-hx}\\
-\partial_x E_z + j\omega\mu H_y &= 0, \label{eq:tm-hy}\\
\partial_x H_y - \partial_y H_x - j\omega\varepsilon E_z &= J_z. \label{eq:tm-ez}
\end{align}
Eliminating $\mathbf{H}$ gives the Helmholtz equation used only as a check,
\begin{equation}
\partial_x^2 E_z + \partial_y^2 E_z + \omega^2\mu\varepsilon E_z = j\omega\mu J_z.
\label{eq:helmholtz-tm}
\end{equation}
The code never forms this scalar operator when $\mu$ varies from face to face.
It keeps the curls.

\paragraph{TEz.}
The unknowns are $H_z$, $E_x$, and $E_y$:
\begin{align}
\partial_x E_y - \partial_y E_x + j\omega\mu H_z &= 0, \label{eq:te-hz}\\
\partial_y H_z - j\omega\varepsilon E_x &= J_x, \label{eq:te-ex}\\
-\partial_x H_z - j\omega\varepsilon E_y &= J_y. \label{eq:te-ey}
\end{align}

\subsection{Constitutive law}
\label{sec:constitutive}

A volumetric law returns $\varepsilon$, $\mu$, and $\sigma$ at the Yee location
of each component. The milestone-1 law is uniform and isotropic: one positive
$\varepsilon$, one positive $\mu$, and one $\sigma\ge 0$. The call also accepts
a frequency and an external state (a temperature map, a stress map, or
anything the caller has already computed). The uniform law ignores both. A
later law may return a different value at every sample, or a complex
$\varepsilon(\omega)$ for FDFD. The time stepper requires real $\varepsilon$
and $\sigma$; a dispersive time-domain medium will add auxiliary variables
rather than a complex coefficient inside the leapfrog.

Surface models are not constitutive samples. PEC removes unknowns. A
resistive sheet adds the conductance of Sec.~\ref{sec:impedance} to one
Ampere cell. A GSTC sheet replaces the magnetic sample on the cut. None of
these is encoded by setting $\varepsilon$ to zero or to infinity.

\subsection{Boundaries in the continuous problem}
\label{sec:continuous-bc}

On a PEC wall the tangential electric field is zero. On a PMC wall the
tangential magnetic field is zero. A periodic wall identifies the two opposite
faces, optionally with a Bloch phase $e^{-j\mathbf{k}_B\cdot\mathbf{L}}$ in a
later slice; milestone~1 uses a phase of one. An open boundary is a PML,
the stretch of Sec.~\ref{sec:pml}.

\section{Discretization}
\label{sec:discretization}

\subsection{Cells, axes, and shapes}
\label{sec:shapes}

The grid has $n_x$ by $n_y$ cells of size $\Delta x$ and $\Delta y$. The domain
is $[0,a]\times[0,b]$ with $a=n_x\Delta x$ and $b=n_y\Delta y$. Stored arrays
use index $[i,j]$ with axis~0 along $x$.

Table~\ref{tab:shapes} is the layout the drivers assume. PEC keeps the
boundary electric samples in the array and sets them to zero. Periodic
identification drops the duplicate far side, so every stored sample is an
unknown and every difference wraps.

\begin{table}[ht]
\centering
\caption{Sample locations. A PEC count of $n{+}1$ includes both boundaries. A
periodic count of $n$ is one period.}
\label{tab:shapes}
\begin{tabular}{@{}llll@{}}
\toprule
Field & Location & PEC shape & Periodic shape\\
\midrule
$E_z$ & $(i\Delta x,\, j\Delta y)$ & $(n_x{+}1,n_y{+}1)$ & $(n_x,n_y)$\\
$H_x$ & $(i\Delta x,\, (j+\tfrac12)\Delta y)$ & $(n_x{+}1,n_y)$ & $(n_x,n_y)$\\
$H_y$ & $((i+\tfrac12)\Delta x,\, j\Delta y)$ & $(n_x,n_y{+}1)$ & $(n_x,n_y)$\\
\midrule
$H_z$ & $((i+\tfrac12)\Delta x,\, (j+\tfrac12)\Delta y)$ & $(n_x,n_y)$ & $(n_x,n_y)$\\
$E_x$ & $((i+\tfrac12)\Delta x,\, j\Delta y)$ & $(n_x,n_y{+}1)$ & $(n_x,n_y)$\\
$E_y$ & $(i\Delta x,\, (j+\tfrac12)\Delta y)$ & $(n_x{+}1,n_y)$ & $(n_x,n_y)$\\
\bottomrule
\end{tabular}
\end{table}

Centered differences of a linear field are exact on this stencil. That is the
sign test for the curls: if $E_z=\alpha x+\beta y$, then the discrete
$\partial_y E_z$ equals $\beta$ and $-\partial_x E_z$ equals $-\alpha$.

\subsection{Discrete curls}
\label{sec:curls}

Collect the free electric samples in a vector $\mathrm{e}$ and the magnetic
samples that the curls touch in a vector $\mathrm{h}$. The two real sparse
matrices are defined so that they match~\eqref{eq:faraday} and~\eqref{eq:ampere}
term by term:
\begin{align}
\mathrm{curl}_E\,\mathrm{e} + j\omega\,\mathrm{diag}(\mu_H)\,\mathrm{h} &= 0,
\label{eq:curl-e}\\
\mathrm{curl}_H\,\mathrm{h} - j\omega\,\mathrm{diag}(\varepsilon_E)\,\mathrm{e}
&= \mathrm{j}.
\label{eq:curl-h}
\end{align}
For TMz, $\mathrm{curl}_E$ stacks $(\partial_y E_z,\,-\partial_x E_z)$ and
$\mathrm{curl}_H$ is $\partial_x H_y-\partial_y H_x$. For TEz,
$\mathrm{curl}_E$ is $\partial_x E_y-\partial_y E_x$ and $\mathrm{curl}_H$
stacks $(\partial_y H_z,\,-\partial_x H_z)$.

On every grid in Table~\ref{tab:shapes} the two matrices are transposes,
\begin{equation}
\mathrm{curl}_H = \mathrm{curl}_E^{\mathsf T}.
\label{eq:adjoint}
\end{equation}
The boundary terms that integration by parts would leave are absent because
PEC electric samples are not unknowns, and periodic faces cancel. The product
$\mathrm{curl}_H\,\mathrm{diag}(\mu^{-1})\,\mathrm{curl}_E$ is therefore
symmetric.

PEC degrees of freedom, TMz:
$E_z$ for $i=1,\ldots,n_x-1$ and $j=1,\ldots,n_y-1$;
$H_x$ for $i=1,\ldots,n_x-1$ and $j=0,\ldots,n_y-1$;
$H_y$ for $i=0,\ldots,n_x-1$ and $j=1,\ldots,n_y-1$.
The magnetic samples on the outer PEC edges do not enter $\mathrm{curl}_H$ at
a free $E_z$. TEz keeps every $H_z$, drops $E_x$ on $y=0$ and $y=b$, and drops
$E_y$ on $x=0$ and $x=a$.

TEz curl-curl on tangential $\mathbf{E}$ has a kernel: any discrete gradient of
a nodal potential that is constant on the PEC boundary has zero circulation.
Those eigenvalues are zero. They are not cavity modes. The physical spectrum
is the positive one.

The 3D cell uses the same partial derivatives on every axis. Sample
locations extend Table~\ref{tab:shapes}: $E_x$ sits at
$((i+\tfrac12)\Delta x,\, j\Delta y,\, k\Delta z)$, and each electric
component is staggered along its own direction. The magnetic components
are staggered along the other two. A PEC wall drops the tangential
electric samples. Magnetic samples that do not touch a free electric
sample are omitted, so~\eqref{eq:adjoint} still holds on a PEC box and on
a periodic box.

\subsection{Leapfrog update}
\label{sec:leapfrog}

Store $\mathbf{E}$ at integer steps and $\mathbf{H}$ at the half-step before
it. One step produces $\mathbf{H}^{n+1/2}$ from $\mathbf{E}^n$, then
$\mathbf{E}^{n+1}$ from $\mathbf{H}^{n+1/2}$. For TMz, with the curls of
Sec.~\ref{sec:curls},
\begin{align}
H_x &\leftarrow H_x - \frac{\Delta t}{\mu_x}\,\partial_y E_z, \label{eq:upd-hx}\\
H_y &\leftarrow H_y + \frac{\Delta t}{\mu_y}\,\partial_x E_z, \label{eq:upd-hy}\\
E_z &\leftarrow C_a E_z + C_b\,(\partial_x H_y - \partial_y H_x). \label{eq:upd-ez}
\end{align}
The coefficients are evaluated at the $E_z$ samples,
\begin{equation}
C_a = \frac{1-\sigma\Delta t/(2\varepsilon)}{1+\sigma\Delta t/(2\varepsilon)},
\qquad
C_b = \frac{\Delta t/\varepsilon}{1+\sigma\Delta t/(2\varepsilon)}.
\label{eq:ca-cb}
\end{equation}
When $\sigma=0$ they are $C_a=1$ and $C_b=\Delta t/\varepsilon$. The same pair
is used on $E_x$ and $E_y$ in TEz, and the magnetic update is
\begin{equation}
H_z \leftarrow H_z - \frac{\Delta t}{\mu_z}(\partial_x E_y - \partial_y E_x).
\label{eq:upd-hz}
\end{equation}
After the electric update, PEC edges are set back to zero.

No matrix is stored. The array curls are the same stencil as
$\mathrm{curl}_E$ and $\mathrm{curl}_H$, including the wrap on a periodic grid.

\subsection{Material sampling and staircase geometry}
\label{sec:sampling}

$\varepsilon$ and $\sigma$ live on electric samples. $\mu$ lives on magnetic
samples. A uniform law fills those arrays with constants. A staircase solid
fills each sample from the coordinate of that sample, so an electric node and
a magnetic node on the same cell can lie on opposite sides of a face. A slab
$y_0\le y<y_1$ is half-open, and a disk $(x-x_c)^2+(y-y_c)^2\le R^2$ is
closed. Later regions overwrite earlier ones. There is no volume-fraction
weight. The assignment is
\begin{equation}
(x,y)\in\mathcal{R}
\quad\Longrightarrow\quad
(\varepsilon,\mu,\sigma)(x,y)=(\varepsilon,\mu,\sigma)_{\mathcal{R}},
\label{eq:staircase}
\end{equation}
The stepper is unchanged, because it already reads one coefficient per sample.
Conformal cuts and sub-cell dielectric smoothing are not part of this scheme.

Tensor $\bar\varepsilon$ is out of scope until the scalar isotropic tests
match the analytic formulas. Off-diagonal samples would sit at locations the
present curls do not read.

\subsection{Bulk Debye and Lorentz media}
\label{sec:dispersion-bulk}

A dispersive law returns a complex $\varepsilon(\omega)$ to FDFD and a real
instantaneous permittivity $\varepsilon_\infty$ to the leapfrog. One electric
Debye pole and one electric Lorentz pole are enough for this slice. With the
$e^{j\omega t}$ convention,
\begin{align}
\varepsilon(\omega)
&= \varepsilon_\infty + \frac{\varepsilon_0\Delta\varepsilon}{1+j\omega\tau},
\label{eq:debye}\\
\tau\partial_t P + P &= \varepsilon_0\Delta\varepsilon\, E,
\label{eq:debye-ode}\\
\varepsilon(\omega)
&= \varepsilon_\infty
+ \frac{\varepsilon_0\omega_p^2}{\omega_0^2-\omega^2+j\gamma\omega},
\label{eq:lorentz-bulk}\\
\partial_t^2 P + \gamma\partial_t P + \omega_0^2 P
&= \varepsilon_0\omega_p^2 E.
\label{eq:lorentz-ode}
\end{align}
Magnetic dispersion is not in this slice. These volume auxiliaries are
independent of the surface oscillators~\eqref{eq:lorentz-e}
and~\eqref{eq:lorentz-m}.

$J_p=\partial_t P$ is added to the impressed current. The coefficients
$C_a$ and $C_b$ stay~\eqref{eq:ca-cb}, built on $\varepsilon_\infty$ and
$\sigma$, and the electric update becomes
\begin{equation}
E \leftarrow C_a E + C_b\,(\nabla\times H - J_p - J).
\label{eq:jp}
\end{equation}
$J=0$ and $J_p=0$ is~\eqref{eq:upd-ez}. Lorentz stores $P$ at integer steps,
advances it explicitly from $E^n$, and sets $J_p^{n+1/2}=(P^{n+1}-P^n)/\Delta t$.
Debye uses the centered average $E^{n+1/2}=(E^{n+1}+E^n)/2$, which couples
$P^{n+1}$ to $E^{n+1}$ and is solved one sample at a time. TEz carries one
auxiliary on $E_x$ and one on $E_y$.

\section{Stability, dispersion, and domain closure}
\label{sec:stability}

\subsection{CFL limit}
\label{sec:cfl}

On a uniform 2D Yee grid the leapfrog is stable for
\begin{equation}
\Delta t \le \Delta t_{\mathrm{CFL}}
= \frac{1}{c\sqrt{\Delta x^{-2}+\Delta y^{-2}}}.
\label{eq:cfl}
\end{equation}
The bound is sharp for the grid mode $k_x=\pi/\Delta x$, $k_y=\pi/\Delta y$.
A run at $0.99\,\Delta t_{\mathrm{CFL}}$ must stay bounded. A run at
$1.01\,\Delta t_{\mathrm{CFL}}$ must grow.

\subsection{Semi-discrete and fully discrete frequencies}
\label{sec:dispersion}

With continuous time and Yee differences, a mode $k_x$, $k_y$ oscillates at
\begin{equation}
\left(\frac{\omega_0}{c}\right)^2
=
\left(\frac{2\sin(k_x\Delta x/2)}{\Delta x}\right)^2
+
\left(\frac{2\sin(k_y\Delta y/2)}{\Delta y}\right)^2.
\label{eq:semi}
\end{equation}
The leapfrog moves that frequency to the root of
\begin{equation}
\left(\frac{\sin(\omega\Delta t/2)}{c\,\Delta t}\right)^2
=
\left(\frac{\sin(k_x\Delta x/2)}{\Delta x}\right)^2
+
\left(\frac{\sin(k_y\Delta y/2)}{\Delta y}\right)^2,
\label{eq:discrete}
\end{equation}
provided the right-hand side is at most $(c\Delta t)^{-2}$. Because
$\arcsin z>z$ on $(0,1)$, this $\omega$ lies slightly above $\omega_0$.
FDFD, having no time difference, must reproduce $\omega_0$.

A rectangular PEC cavity of size $a\times b$ has $k_x=m\pi/a$ and
$k_y=n\pi/b$. TMz needs $m\ge 1$ and $n\ge 1$, because $E_z$ vanishes on the
whole boundary. TEz also allows $m=0$ or $n=0$, but not both. The mode shapes
consistent with the stencil are
\begin{align}
E_z &= \sin(m\pi x/a)\sin(n\pi y/b), \label{eq:mode-tm}\\
H_z &= \cos(m\pi x/a)\cos(n\pi y/b). \label{eq:mode-te}
\end{align}
A rectangular box $a\times b\times c$ adds the term
$(2\sin(k_z\Delta z/2)/\Delta z)^2$ to~\eqref{eq:semi}, with
$k_z=p\pi/c$. The indices are non-negative and at least two of them are
positive. A mode with exactly one index equal to zero is a single
Cartesian component: $(m,n,0)$ is $E_z=\sin(m\pi x/a)\sin(n\pi y/b)$,
$(m,0,p)$ is $E_y$, and $(0,n,p)$ is $E_x$. These three are exact
eigenmodes of the centered difference. A mode with three positive indices
has two polarizations linked by $k_x A_x+k_y A_y+k_z A_z=0$.

A pure spatial eigenmode sampled on three successive steps obeys the scalar
leapfrog recurrence. Those three numbers determine $\omega$ without a
Fourier transform.

\subsection{Lossless invariant}
\label{sec:invariant}

$\mathbf{E}$ and $\mathbf{H}$ are never stored at the same time. The sum of
their sample energies therefore oscillates by a noticeable fraction of itself
near the CFL limit, even though the scheme is lossless. The quadratic form
that the update conserves, when $\sigma=0$, is
\begin{equation}
W^n
=
\frac12\sum \varepsilon\,(E^n)^2\,\Delta x\Delta y
+
\frac12\sum \mu\, H^{n+1/2} H^{n-1/2}\,\Delta x\Delta y,
\label{eq:invariant}
\end{equation}
with the sums over the stored samples of each component. The derivation is the
adjoint relation~\eqref{eq:adjoint}: the work done on $\mathbf{E}$ by
$\mathrm{curl}_H\mathbf{H}$ cancels the work done on $\mathbf{H}$ by
$\mathrm{curl}_E\mathbf{E}$. Conductivity removes the cancellation. Above the
CFL limit $W^n$ can stay constant while the ordinary field energy grows,
because $W^n$ is then indefinite. Boundedness is checked with $W^n$. Growth
past the CFL limit is checked with the ordinary field energy.

\subsection{PEC, PMC, and periodic walls}
\label{sec:walls}

PEC for TMz is $E_z=0$ on the outer node lines. PEC for TEz is $E_x=0$ on
$y\in\{0,b\}$ and $E_y=0$ on $x\in\{0,a\}$. PMC is the dual statement,
tangential $\mathbf{H}=0$, which is a Neumann condition on the remaining
field. PMC is part of the surface stack and is not in the milestone-1
unknown set.

Periodic differences replace $i-1$ at $i=0$ by $i=n_x-1$, and likewise in $y$.
A Bloch phase multiplies the wrapped sample by $e^{\pm j k_{B,x} a}$. The
phase is unity until the metasurface unit-cell slice.

\subsection{PML}
\label{sec:pml}

An open domain is closed by a complex stretch of the normal coordinate,
following the convolutional PML. In the frequency domain, for each axis
$w\in\{x,y\}$,
\begin{equation}
s_w = \kappa_w + \frac{\sigma_w}{\alpha_w + j\omega\varepsilon_0},
\label{eq:stretch}
\end{equation}
and every $\partial_w$ in the curls is replaced by $s_w^{-1}\partial_w$. The
reciprocal that the auxiliary equation integrates is
\begin{equation}
\frac{1}{s_w}
=
\frac{1}{\kappa_w}
-
\frac{\sigma_w/(\kappa_w^2\varepsilon_0)}
{j\omega + \alpha_w/\varepsilon_0 + \sigma_w/(\kappa_w\varepsilon_0)}.
\label{eq:stretch-inv}
\end{equation}
Inside a PML slab of thickness $d$, with $\rho$ the depth measured from the
inner face and $m$ typically 3 or 4,
\begin{equation}
\sigma_w(\rho) = \sigma_{\max}(\rho/d)^m,
\qquad
\sigma_{\max} = -\frac{(m+1)\ln R(0)}{2\eta_0 d},
\label{eq:sigma-max}
\end{equation}
with a target reflection $R(0)$ of order $10^{-6}$ to $10^{-8}$. $\kappa_w$
grows from 1 at the inner face to a modest $\kappa_{\max}$, and $\alpha_w$
falls from a small positive value at the inner face to zero at the outer PEC
backing. The time-domain scheme adds one auxiliary per stretched derivative
and integrates~\eqref{eq:stretch-inv} with the same $\Delta t$ as the
leapfrog. \texttt{operators.pml} is that stretch, on a PEC-backed band, in
both drivers. A solve with no PML object is the milestone-1 stepper. The
regression compares a line current with the outgoing $H_0^{(2)}$ field on
the unstretched samples.

\section{Linear algebra and thin sheets}
\label{sec:algebra}

\subsection{Driven FDFD system}
\label{sec:fdfd}

Eliminating $\mathrm{h}$ from~\eqref{eq:curl-e} and~\eqref{eq:curl-h} gives
the system that is factored,
\begin{equation}
\mathrm{curl}_H\,\mathrm{diag}\!\left(\frac{1}{j\omega\mu_H}\right)\mathrm{curl}_E\,\mathrm{e}
+ j\omega\,\mathrm{diag}(\varepsilon_E)\,\mathrm{e}
= -\mathrm{j}.
\label{eq:fdfd}
\end{equation}
The unknown is the free electric vector. PEC rows are omitted rather than
penalized. The matrix is complex symmetric when $\varepsilon$ and $\mu$ are
scalar, and it is stored sparse. Milestone~1 solves it with a sparse direct
factorization. Two-dimensional grids of the size used in the tests factor in
negligible time. Iterative solves (GMRES or BiCGSTAB, with a preconditioner)
wait until a grid is large enough that the factor is the expensive part.

At a fixed real $\omega$ away from the cavity spectrum the matrix is
nonsingular. A manufactured solution builds $\mathrm{j}$ by applying the
left-hand side to a chosen $\mathrm{e}$ and asks the solver for that
$\mathrm{e}$ back.

\subsection{Cavity eigenproblem}
\label{sec:eigen}

For lossless media and $\mathrm{j}=0$, multiplying~\eqref{eq:fdfd} through by
$j\omega$ produces the generalized eigenproblem
\begin{equation}
K\mathrm{e} = \omega^2 M\mathrm{e},
\qquad
K = \mathrm{curl}_H\,\mathrm{diag}(\mu_H^{-1})\,\mathrm{curl}_E,
\qquad
M = \mathrm{diag}(\varepsilon_E).
\label{eq:eigen}
\end{equation}
$K$ is real symmetric and, on a TMz PEC cavity, positive definite. On a TEz
PEC cavity it is positive semi-definite with the gradient kernel of
Sec.~\ref{sec:curls}. The positive eigenvalues are the semi-discrete
frequencies~\eqref{eq:semi}. Conductivity does not enter $K$ and $M$. A
complex permittivity belongs in the driven system~\eqref{eq:fdfd}, not in
this real eigenproblem.

\subsection{Resistive sheet}
\label{sec:impedance}

A zero-thickness electric sheet with surface impedance $Z_s$ (ohms per
square) carries
\begin{equation}
\hat{\mathbf{n}}\times(\mathbf{H}^+-\mathbf{H}^-) = \mathbf{E}_{\mathrm{av}}/Z_s.
\label{eq:zs}
\end{equation}
The normal points from the $-$ side to the $+$ side. For a sheet of normal
$\hat{\mathbf{y}}$ in TMz this is $-(H_x^+-H_x^-)=E_z/Z_s$.

The same sheet is the pure-electric GSTC of Sec.~\ref{sec:gstc} with
\begin{equation}
\chi_{ee} = \frac{\eta_0}{j k Z_s} = \frac{1}{j\omega\varepsilon_0 Z_s},
\label{eq:chi-from-z}
\end{equation}
and $\chi_{mm}=0$. MaxwellMOM's \texttt{chi\_ee\_from\_Zs} is this map. In the
time domain a real $Z_s$ is a shunt conductance, not a constant
susceptibility: $J_z=\sigma_s E_z$ with $\sigma_s=1/Z_s$. Placing that current
in the Ampere cell of height $\Delta y$ that contains the sheet adds
\begin{equation}
\sigma_{\mathrm{eff}} = \sigma + \frac{1}{Z_s\Delta y}
\label{eq:sigma-sheet}
\end{equation}
to the $C_a$, $C_b$ formula~\eqref{eq:ca-cb} on the sheet's electric samples.
That is the milestone-4 update. It is the $\chi_{mm}=0$ reduction of the
symmetric cell, specialized to $\chi_{ee}\propto 1/(j\omega)$.

\subsection{Complex permittivity and a Dirichlet trace}
\label{sec:dirichlet}

When $\varepsilon$ and $\sigma$ are real, the diagonal of~\eqref{eq:fdfd} uses
\begin{equation}
\varepsilon_{\mathrm{FDFD}} = \varepsilon + \frac{\sigma}{j\omega}.
\label{eq:eps-sigma}
\end{equation}
A law that returns a complex $\varepsilon(\omega)$ has already included its
loss. That array replaces~\eqref{eq:eps-sigma} and the accompanying $\sigma$
is zero, so a pole is not counted twice. The real eigenproblem~\eqref{eq:eigen}
still requires real positive $\varepsilon$ and $\mu$ and $\sigma=0$.

PEC electric samples stay out of the unknown vector. A non-zero trace $e_D$,
equal to the prescribed boundary values and zero on the free samples, moves
onto the right-hand side,
\begin{equation}
A_{ff}\,\mathrm{e}_f = -\mathrm{j} - (A\mathrm{e}_D)_f.
\label{eq:dirichlet}
\end{equation}
The free rows of $A\mathrm{e}_D$ are the array form of the left-hand side
of~\eqref{eq:fdfd}. A zero trace is the milestone-1 solve. This is how the
slab and cylinder comparisons are posed. It is not a PML.

An embedded PEC uses the same omission on every electric sample whose
coordinate lies inside the object. An embedded PMC omits the magnetic
samples in its region, so the tangential $H$ on a staircase face is zero
and the magnetic samples of a solid interior leave the unknown vector.
Both are masks on the Yee locations. Setting $\varepsilon$ to infinity, or
$\mu$ to zero, is a different model. The open PEC cylinder puts
$e_D=-E_{\mathrm{inc}}$ on the object and $e_D=0$ on the outer wall, and
the PML absorbs the scattered field. Inside the object the total electric
field is zero.

An open dielectric cylinder keeps that PML and the homogeneous outer wall.
It does not remove interior samples. The staircase permittivity stays on
the diagonal of~\eqref{eq:fdfd}, and the incident wave enters as the
impressed current $J=j\omega(\varepsilon-\varepsilon_0)E_{\mathrm{inc}}$
on the electric samples. The current is zero wherever
$\varepsilon=\varepsilon_0$. A free sample of the total field is
$E_s+E_{\mathrm{inc}}$. Inside the disk the total field is the Bessel
series of Sec.~\ref{sec:mie2d}.

\section{Validation}
\label{sec:validation}

Milestone~1 is a numerical check that the curls, the leapfrog, and the sparse
system implement Secs.~\ref{sec:discretization} and~\ref{sec:stability}. The
acceptance tests and the bars they use are:

\begin{center}
\small
\begin{tabular}{@{}>{\raggedright\arraybackslash}p{5.0cm}>{\raggedright\arraybackslash}p{3.5cm}>{\raggedright\arraybackslash}p{6.2cm}@{}}
\toprule
Check & Bar & What is compared\\
\midrule
Linear TMz and TEz curls & $10^{-12}$ absolute & Centered differences of a linear field\\
$\mathrm{curl}_H=\mathrm{curl}_E^{\mathsf T}$ & $10^{-12}$ absolute & All four grids in Table~\ref{tab:shapes}\\
Matrix versus array curl & $10^{-12}$ relative & Same four grids, PEC boundaries held at zero\\
TMz cavity $(1,1)$ & $10^{-10}$ relative & Smallest FDFD eigenvalue versus~\eqref{eq:semi}\\
TEz cavities $(1,0)$ and $(1,1)$ & $10^{-10}$ relative & Positive FDFD eigenvalues versus~\eqref{eq:semi}\\
Manufactured FDFD solve & $10^{-10}$ relative & Recovered $\mathrm{e}$ on TMz and TEz\\
Leapfrog invariant at $0.99\,\Delta t_{\mathrm{CFL}}$ & 1 percent & $W^n$ over $2\cdot 10^3$ steps, random TMz data\\
Field energy at $1.01\,\Delta t_{\mathrm{CFL}}$ & growth past $10\times$ & Same data, ordinary sample energy\\
TMz and TEz modal $\omega$ & $10^{-6}$ relative & Three leapfrog samples versus~\eqref{eq:discrete}\\
Periodic standing wave $(k_x,k_y)=(2\pi/a,\,0)$ & $10^{-6}$ relative & Same recurrence versus~\eqref{eq:discrete}\\
\bottomrule
\end{tabular}
\end{center}

The modal and standing-wave bars in the project plan were $10^{-3}$, which is
a fit tolerance. A pure eigenmode has no fit error, so the tests use
$10^{-6}$.

Milestone~2 adds the volumetric checks. The scattering rows compare the Yee
solution, with the trace~\eqref{eq:dirichlet} taken from the formulas below,
to those formulas. The gap is dispersion, and for the disk the staircase.
A slab face used in a convergence pair is kept off the electric nodes of
both grids, so~\eqref{eq:staircase} does not assign the face to a sample.
The comparison wavelength is kept off the TEz eigenfrequencies of that
closed box: $\lambda=40\Delta$ on the $8\times 80$ domain is one of them,
and the Dirichlet problem is then singular.

\begin{center}
\small
\begin{tabular}{@{}>{\raggedright\arraybackslash}p{5.0cm}>{\raggedright\arraybackslash}p{3.5cm}>{\raggedright\arraybackslash}p{6.2cm}@{}}
\toprule
Check & Bar & What is compared\\
\midrule
Fresnel face continuity & $10^{-12}$ absolute & Tangential $E$ and $H$, both polarizations\\
Mie cylinder continuity & $10^{-8}$ absolute & Tangential field and derivative at $\rho=a$\\
PEC cylinder coefficient & $10^{-12}$ relative & $a_n$ against $-(-j)^n J_n/H_n^{(2)}$\\
Manufactured $\sigma>0$ and $\varepsilon(\mathbf{r})$ & $10^{-10}$ relative & Recovered $\mathrm{e}$, both polarizations\\
Uniform $\varepsilon_r=4$ cavity & $10^{-10}$ / $10^{-6}$ & FDFD versus~\eqref{eq:semi}; leapfrog versus~\eqref{eq:discrete}\\
Slab-loaded TMz cavity & $10^{-2}$ and $10^{-6}$ & Continuous root, then leapfrog against that $\omega_0$\\
Lossless staircase invariant & 1 percent & $W^n$ over 500 steps at half the CFL step\\
Fresnel slab, both polarizations & $2\cdot 10^{-2}$ relative & Free samples; halving $\Delta$ drops the error by $2.5\times$\\
Dielectric cylinder, both polarizations & $2\cdot 10^{-2}$ relative & Samples $2\Delta$ off the circle and the wall\\
Debye and Lorentz, one period & $3\cdot 10^{-3}$ relative & Leapfrog against the complex-$\varepsilon$ phasor\\
\bottomrule
\end{tabular}
\end{center}

The convolutional PML is checked apart from those Dirichlet rows. The
profile row is the grading of Sec.~\ref{sec:pml}. The line-current rows
compare $E_z$ of a unit line current with the outgoing $H_0^{(2)}$ field
on the unstretched samples at least $\lambda/2$ from the source. Halving
the spacing keeps the same physical band. The bare-box row uses that same
window with the thicknesses set to zero. The leapfrog row is a one-period
transform, after the transient, against the FDFD phasor at $48$ steps per
period.

\begin{center}
\small
\begin{tabular}{@{}>{\raggedright\arraybackslash}p{5.0cm}>{\raggedright\arraybackslash}p{3.5cm}>{\raggedright\arraybackslash}p{6.2cm}@{}}
\toprule
Check & Bar & What is compared\\
\midrule
CFS profile & $10^{-12}$ relative & $\sigma$, $\kappa$, $\alpha$ on the node and midpoint lines\\
Stretched curls, both polarizations & $10^{-12}$ relative & Sparse pair versus the array curls\\
Zero-thickness PML step & exact match & Unstretched leapfrog, both polarizations\\
Line current, coarse band & $4\cdot 10^{-2}$ relative & Interior $E_z$ versus $H_0^{(2)}$\\
Line current, refined band & $2.5\times$ smaller & Same physical band, half the spacing\\
Bare PEC window & above $0.5$ relative & Those interior samples, with no stretch\\
Leapfrog phasor, both polarizations & $1.5\cdot 10^{-2}$ relative & $48$ steps per period versus FDFD\\
\bottomrule
\end{tabular}
\end{center}

Embedded conductors are checked on the reduced curl. A PEC band that fills
$y\ge b_{\mathrm{short}}$ on a taller grid leaves the spectrum of a PEC
cavity of height $b_{\mathrm{short}}$. A PMC band whose lower edge sits on
an $H_x$ row, $y=(J+\tfrac12)\Delta y$, leaves the TMz quarter-wave
$k_y=(n+\tfrac12)\pi/y_{\mathrm{face}}$. Decoupled electric samples above
that cut are zero eigenvalues and are dropped before the comparison. The
open rows use a $64\times 64$ grid with $\Delta=1$, $a=8$, eight PML
cells, and $ka=1$. The series takes $k=\omega\sqrt{\mu_0\varepsilon_0}$.
The PEC window is the unstretched samples with $\rho\ge a+2\Delta$. The
dielectric rows use that grid with $n=1.5$. Their window is the
unstretched samples with $|\rho-a|\ge 2\Delta$, inside and outside the
disk, and the drive is the contrast current of Sec.~\ref{sec:dirichlet}.

\begin{center}
\small
\begin{tabular}{@{}>{\raggedright\arraybackslash}p{5.0cm}>{\raggedright\arraybackslash}p{3.5cm}>{\raggedright\arraybackslash}p{6.2cm}@{}}
\toprule
Check & Bar & What is compared\\
\midrule
PEC band, both polarizations & $10^{-10}$ relative & Shortened cavity versus~\eqref{eq:semi}\\
PMC quarter-wave, TMz & $10^{-10}$ relative & Lowest positive eigenvalue, $k_y=\pi/(2y_{\mathrm{face}})$\\
Open PEC cylinder, TMz & $8\cdot 10^{-2}$ relative & $E_z$, $\rho\ge a+2\Delta$\\
Open PEC cylinder, TEz & $10^{-1}$ relative & $E_x$ and $E_y$, same window\\
Open dielectric cylinder, TMz & $8\cdot 10^{-3}$ relative & $E_z$, $|\rho-a|\ge 2\Delta$\\
Open dielectric cylinder, TEz & $1.3\cdot 10^{-2}$ relative & $E_x$ and $E_y$, same window\\
Resistive sheet, TMz & $6\cdot 10^{-4}$ relative & $E_z$, $8\times 80$, $\Delta=1$, $Z_s=\eta_0$, $\lambda=46$\\
Resistive sheet, TEz & $1.2\cdot 10^{-2}$ relative & $E_x$ and $E_y$, same grid\\
Symmetric sheet, TMz & $4.2\cdot 10^{-2}$ relative & $E_z$, same box, face $y=40.5$, $\chi_{ee}=0.5$, $\chi_{mm}=0.25$\\
Symmetric sheet, TEz & $1.9\cdot 10^{-1}$ relative & $E_x$ and $E_y$, same grid\\
\bottomrule
\end{tabular}
\end{center}

The resistive-sheet rows are a closed PEC box with zero impressed current
and the analytic trace on the wall. The sheet is the
conductance~\eqref{eq:sigma-sheet} on one Ampere cell, at $y=40$. The
comparison is the packed free electric vector. The analytic coefficients
at $Z_s=100\,\Omega$ and $k=2\pi/0.3$ match one stored MaxwellMOM value of
that sheet. The symmetric-sheet rows use the same box with the sheet on
the magnetic face $y=40.5$, $\chi_{ee}=0.5$, and $\chi_{mm}=0.25$. The
comparison is again the packed free electric vector. A cell of $\Delta=1/3$
on that same sheet drops both errors by more than half. The coefficients
at $k=2\pi/0.3$ for this $\chi$ match one stored MaxwellMOM value.

The 3D curls use the partial derivatives of Sec.~\ref{sec:curls} on every
axis. The cavity row compares the Rayleigh quotient of a known mode with
\eqref{eq:semi} after the term $(2\sin(k_z\Delta z/2)/\Delta z)^2$ is
added, with $k_z=p\pi/c$. The box is $8\times 6\times 5$ cells at
$\Delta=1$, and the modes are $(1,1,0)$, $(1,0,1)$, and $(0,1,1)$. Each
has one index equal to zero and a single Cartesian component. The
manufactured row recovers a random electric vector on a $6\times 5\times 4$
box with spacings $0.5$, $0.4$, and $0.3$.

\begin{center}
\small
\begin{tabular}{@{}>{\raggedright\arraybackslash}p{5.0cm}>{\raggedright\arraybackslash}p{3.5cm}>{\raggedright\arraybackslash}p{6.2cm}@{}}
\toprule
Check & Bar & What is compared\\
\midrule
3D linear curls & $10^{-12}$ absolute & Centered differences of a linear field\\
3D $\mathrm{curl}_H=\mathrm{curl}_E^{\mathsf T}$ & $10^{-12}$ absolute & PEC and periodic boxes\\
3D matrix versus array curl & $10^{-12}$ relative & Same boxes, PEC boundaries held at zero\\
3D PEC cavity, three pure modes & $10^{-10}$ relative & $8\times 6\times 5$ box, versus~\eqref{eq:semi} plus $z$\\
3D manufactured FDFD solve & $10^{-10}$ relative & Recovered $\mathrm{e}$ on a $6\times 5\times 4$ box\\
3D stretched curls & $10^{-12}$ relative & Matrix versus array curl, $6\times 5\times 4$ box\\
3D Mie surface & $10^{-8}$ absolute & Tangential $E,H$ and normal $D,B$ at $\rho=a$\\
Open dielectric sphere & $1.8\cdot 10^{-2}$ relative & $E$, $|\rho-a|\ge 2\Delta$, $18^3$ box, PML of 3 cells\\
\bottomrule
\end{tabular}
\end{center}

The stretched-curl row multiplies each partial derivative by $1/s_w$ and
compares the sparse matrix with the array curl. The surface row evaluates
the sphere series from both sides of $\rho=a$. Tangential $E$ and $H$
match, as do the normal components of $D$ and $B$. That series is the
outgoing spherical wave $h_n^{(2)}$. Its coefficients are not written out
in this note. The open sphere is a contrast-current solve: the staircase
ball stays in the permittivity, the unknown is the scattered field, and
the comparison keeps the electric samples outside the PML and at least
$2\Delta$ off the surface.

\subsection{Normal-incidence Fresnel slab}
\label{sec:fresnel}

A slab $y_1\le y\le y_2$ of thickness $d$, relative permittivity
$\varepsilon_r$, and $\mu=\mu_0$ sits in vacuum. The incident wave travels
toward $+y$ and equals $1$ at the front face. With $n=\sqrt{\varepsilon_r}$,
$k=\omega/c$, and $\delta=nkd$, a single-interface reflection $r$ produces
\begin{equation}
R = \frac{r(1-e^{-2j\delta})}{1-r^2 e^{-2j\delta}},
\qquad
\tau = \frac{(1-r^2)\,e^{-j\delta}}{1-r^2 e^{-2j\delta}}.
\label{eq:fresnel-R}
\end{equation}
For TMz the unknown is $E_z$ and $r=(1-n)/(1+n)$. The field is
\begin{equation}
E_z(y) =
\begin{cases}
e^{-jk(y-y_1)} + R\, e^{jk(y-y_1)}, & y\le y_1,\\[0.3em]
A\, e^{-jnk(y-y_1)} + C\, e^{jnk(y-y_1)}, & y_1\le y\le y_2,\\[0.3em]
\tau\, e^{-jk(y-y_2)}, & y\ge y_2,
\end{cases}
\label{eq:fresnel-field}
\end{equation}
with $A=\tfrac12\big(1+R+(1-R)/n\big)$ and
$C=\tfrac12\big(1+R-(1-R)/n\big)$. Then
$H_x=-\partial_y E_z/(j\omega\mu_0)$: a forward vacuum wave has
$H_x=E_z/\eta$ and a backward wave has $H_x=-E_z/\eta$, where
$\eta=\sqrt{\mu_0/\varepsilon_0}$. Inside, the forward part of $H_x$ carries
an extra factor of $n$.

TEz is a different scalar. Tangential $H$ reflects with
$r_H=(n-1)/(n+1)=-r$, and~\eqref{eq:fresnel-R} is evaluated at $r_H$.
Transmission depends on $r^2$, so $\tau$ is unchanged and $R_H=-R$.
Interior amplitudes satisfy $A_H+C_H=1+R_H$ and $A_H-C_H=n(1-R_H)$,
\begin{equation}
H_z(y) =
\begin{cases}
e^{-jk(y-y_1)} + R_H\, e^{jk(y-y_1)}, & y\le y_1,\\[0.3em]
A_H\, e^{-jnk(y-y_1)} + C_H\, e^{jnk(y-y_1)}, & y_1\le y\le y_2,\\[0.3em]
\tau\, e^{-jk(y-y_2)}, & y\ge y_2,
\end{cases}
\label{eq:fresnel-te}
\end{equation}
with $A_H=\tfrac12\big(1+R_H+n(1-R_H)\big)$ and
$C_H=\tfrac12\big(1+R_H-n(1-R_H)\big)$. The lift uses
$E_x=\partial_y H_z/(j\omega\varepsilon)$ and $E_y=0$. Outside, a forward
wave has $E_x=-\eta H_z$; inside, that factor is $\eta/n$.
When $\varepsilon_r=1$, both reflections vanish and
$\tau=e^{-jkd}$, so each field equals the incident wave. Continuity of
tangential $E$ and $H$ at both faces is the check on these closed forms.

\subsection{Dielectric cylinder}
\label{sec:mie2d}

The outgoing cylindrical wave is $H_n^{(2)}$, the same convention as
MaxwellMOM's two-dimensional PEC cylinder. The incident field
$e^{-jkx}$ has the Jacobi--Anger expansion
\begin{equation}
e^{-jkx} = \sum_{n=-\infty}^{\infty} (-j)^n J_n(k\rho)\, e^{jn\varphi}.
\label{eq:mie-inc}
\end{equation}
Our TMz unknown is $E_z$. For a non-magnetic cylinder of radius $a$ and
index $m=\sqrt{\varepsilon_r}$, with $\alpha=ka$ and $\beta=mka$,
\begin{equation}
a_n = (-j)^n
\frac{m J_n'(\beta)\, J_n(\alpha) - J_n(\beta)\, J_n'(\alpha)}
{J_n(\beta)\, H_n^{(2)\prime}(\alpha) - m J_n'(\beta)\, H_n^{(2)}(\alpha)}.
\label{eq:mie-tm}
\end{equation}
The interior coefficient follows from continuity of $E_z$ at $\rho=a$.
Outside, the scattered field is $a_n H_n^{(2)}(k\rho)\, e^{jn\varphi}$.
Inside, it is $b_n J_n(mk\rho)\, e^{jn\varphi}$.

TEz is the $H_z$ problem. Tangential $E$ brings $1/\varepsilon$, so the
contrast multiplies the undifferentiated interior Bessel function,
\begin{equation}
a_n = (-j)^n
\frac{m J_n(\beta)\, J_n'(\alpha) - J_n'(\beta)\, J_n(\alpha)}
{J_n'(\beta)\, H_n^{(2)}(\alpha) - m J_n(\beta)\, H_n^{(2)\prime}(\alpha)}.
\label{eq:mie-te}
\end{equation}
The PEC reduction used as a regression is the separate coefficient
$a_n=-(-j)^n J_n(\alpha)/H_n^{(2)}(\alpha)$, not the limit
$\varepsilon_r\to\infty$ of~\eqref{eq:mie-tm}. The TEz form, from
$E_\varphi=0$, is
$a_n=-(-j)^n J_n'(\alpha)/H_n^{(2)\prime}(\alpha)$.

The open PEC cylinder is the scattered-field pair in the conductor table
above. The open dielectric cylinder is the contrast-current pair in that
same table. The open dielectric sphere is the contrast-current row in the
3D table. The resistive sheet and the symmetric sheet are the last rows
in the table above the 3D curls. Later slices add comparisons that this
note already fixes the formulas for, and that the code does not yet run:

\begin{itemize}
\item a dipole, a Gaussian beam, a near-to-far integral, and a Bloch cell.
\end{itemize}

Prescribed external maps are checked by confirming that the samples the
stepper reads are the samples the law was given. They are not a coupled
thermal or mechanical solve.

\section{GSTC metasurfaces}
\label{sec:gstc}

\subsection{Jump conditions}
\label{sec:jumps}

A zero-thickness sheet of normal $\hat{\mathbf{n}}$, pointing from the $-$
side to the $+$ side, replaces the tangential continuity conditions by
\begin{align}
\hat{\mathbf{n}}\times(\mathbf{H}^+-\mathbf{H}^-) &= \partial_t\mathbf{P}_s,
\label{eq:gstc-h}\\
(\mathbf{E}^+-\mathbf{E}^-)\times\hat{\mathbf{n}} &= \mu_0\,\partial_t\mathbf{M}_s.
\label{eq:gstc-e}
\end{align}
For a monoanisotropic sheet the surface polarizations are
\begin{equation}
\mathbf{P}_s = \varepsilon_0\,\chi_{ee}\cdot\mathbf{E}_{\mathrm{av}},
\qquad
\mathbf{M}_s = \chi_{mm}\cdot\mathbf{H}_{\mathrm{av}},
\label{eq:pol}
\end{equation}
with the average of the two one-sided tangential fields. In the frequency
domain $\partial_t$ is $j\omega$. These are the same jumps as in the symmetric
Yee-cell scheme of Smy, Stewart, and Gupta (arXiv:1706.10136) and the same
$e^{j\omega t}$ normalization as MaxwellMOM's \texttt{rt\_from\_chi}: there
$J=j\omega\varepsilon_0\chi_{ee}E_{\mathrm{av}}$ and
$M=j\omega\mu_0\chi_{mm}H_{\mathrm{av}}$.

For TMz and a sheet of normal $\hat{\mathbf{y}}$, \eqref{eq:gstc-h}
and~\eqref{eq:gstc-e} become
\begin{align}
-(H_x^+ - H_x^-) &= \partial_t P_z, \label{eq:jump-hx}\\
E_z^+ - E_z^- &= -\mu_0\,\partial_t M_x. \label{eq:jump-ez}
\end{align}
$H_y$ is the normal field. It is not split.

Bianisotropic coupling, when it is added, extends~\eqref{eq:pol} to
\begin{equation}
\mathbf{P}_s = \varepsilon_0\chi_{ee}\cdot\mathbf{E}_{\mathrm{av}}
+ c_0^{-1}\chi_{em}\cdot\mathbf{H}_{\mathrm{av}},
\qquad
\mathbf{M}_s = \chi_{mm}\cdot\mathbf{H}_{\mathrm{av}}
+ \eta_0^{-1}\chi_{me}\cdot\mathbf{E}_{\mathrm{av}}.
\label{eq:bianiso}
\end{equation}
The reflective lock used in MaxwellMOM is a particular $\chi_{em}$ inside this
extension. It is not part of the first sheet implementation.

\subsection{Reflection and transmission}
\label{sec:rt}

Take a uniform sheet at $y=0$, host medium $(\varepsilon,\mu)$ on both sides,
and a TMz wave at angle $\theta$ from the normal. The reflected and
transmitted amplitudes are defined by $E_z^+=T E_0$ and $E_z^-=E_0(1+R)$ at
the sheet, with the incident wave coming from $y<0$. Substituting the plane
waves into~\eqref{eq:jump-hx} and~\eqref{eq:jump-ez}, with constant
susceptibilities, gives
\begin{equation}
R+T = \frac{1-\tilde\alpha}{1+\tilde\alpha},
\qquad
T-R = \frac{1-\tilde\beta}{1+\tilde\beta},
\label{eq:rt-system}
\end{equation}
where $\alpha_0 = jk\chi_{ee}/2$, $\beta_0 = jk\chi_{mm}/2$, $k=\omega/c$, and
\begin{equation}
\tilde\alpha = \alpha_0/\cos\theta,
\qquad
\tilde\beta = \beta_0\cos\theta.
\label{eq:te-angle}
\end{equation}
Hence
\begin{align}
R &= \frac12\left(
\frac{1-\tilde\alpha}{1+\tilde\alpha}
-
\frac{1-\tilde\beta}{1+\tilde\beta}
\right), \label{eq:R}\\
T &= \frac12\left(
\frac{1-\tilde\alpha}{1+\tilde\alpha}
+
\frac{1-\tilde\beta}{1+\tilde\beta}
\right). \label{eq:T}
\end{align}
At $\theta=0$ the tildes drop off and~\eqref{eq:R}--\eqref{eq:T} are
MaxwellMOM's normal-incidence formula. TMz with the sheet normal
$\hat{\mathbf{y}}$ and the plane of incidence the $xy$-plane is $s$
polarization, which that code calls TE.

TEz on the same sheet is $p$ polarization (TM in MaxwellMOM). The cosines
trade places,
\begin{equation}
\tilde\alpha = \alpha_0\cos\theta,
\qquad
\tilde\beta = \beta_0/\cos\theta,
\label{eq:tm-angle}
\end{equation}
and $R$, $T$ keep the form~\eqref{eq:R}--\eqref{eq:T}. A normal magnetic
susceptibility $\chi_{mm}^{nn}$ enters only the $s$/TE electric coefficient,
\begin{equation}
\tilde\alpha \leftarrow \tilde\alpha + \frac{jk\sin^2\theta\,\chi_{mm}^{nn}}{2\cos\theta},
\label{eq:chi-nn}
\end{equation}
which is the correction MaxwellMOM applies for a flat sheet. The coded TMz
jump keeps $H_y$ unsplit and adds $\chi_{mm}^{nn}\partial_x H_{y,\mathrm{av}}$
to~\eqref{eq:sc-h}. Equation~\eqref{eq:chi-nn} is the plane-wave reduction of
that term. TEz does not carry it.

\subsection{Symmetric cell}
\label{sec:symmetric-cell}

The sheet sits halfway between electric nodes $E_z^-$ at index $j$ and
$E_z^+$ at index $j+1$. The magnetic sample that would have occupied that face
is replaced by a pair $H_x^-$, $H_x^+$ on the two sides of a zero-thickness
cut. The averages and jumps in~\eqref{eq:jump-hx} and~\eqref{eq:jump-ez} use
this pair and the two electric nodes. Bulk Faraday updates for $H_y$, and for
the $H_x$ faces that are not the cut, are unchanged: their stencils do not
cross the sheet.

In the frequency domain the cut and the two Ampere cells that touch it are
the local system
\begin{align}
H_x^+ - H_x^- + j\omega\varepsilon_0\chi_{ee}\frac{E_z^-+E_z^+}{2} &= 0,
\label{eq:sc-h}\\
E_z^+ - E_z^- + j\omega\mu_0\chi_{mm}\frac{H_x^-+H_x^+}{2} &= 0,
\label{eq:sc-e}\\
\frac{H_y|_{i,j}-H_y|_{i-1,j}}{\Delta x}
- \frac{H_x^- - H_x^{\mathrm{below}}}{\Delta y}
- j\omega\varepsilon E_z^- &= J_z^-,
\label{eq:sc-amp-m}\\
\frac{H_y|_{i,j+1}-H_y|_{i-1,j+1}}{\Delta x}
- \frac{H_x^{\mathrm{above}} - H_x^+}{\Delta y}
- j\omega\varepsilon E_z^+ &= J_z^+.
\label{eq:sc-amp-p}
\end{align}
$H_x^{\mathrm{below}}$ and $H_x^{\mathrm{above}}$ are the ordinary Yee samples
one cell off the cut. They are known from the bulk Faraday update when the
cell is marched in time, and they are ordinary unknowns when the same rows are
assembled into~\eqref{eq:fdfd}. Equations~\eqref{eq:sc-h} and~\eqref{eq:sc-e}
replace the Faraday row that used to connect $E_z^-$ to $E_z^+$ across one
$\Delta y$.

For a non-dispersive real $\chi$ the time-domain statement is the same pair
with $j\omega$ replaced by $\partial_t$ acting on $\mathbf{P}_s$ and
$\mathbf{M}_s$. The four surface unknowns at a given step are coupled, so the
cell is a $4\times 4$ (or a $2\times 2$ after eliminating $H_x^\pm$) solve,
while the rest of the grid stays an explicit leapfrog. That local solve is
the one described for the symmetric cell in arXiv:1706.10136.

Two other placements are on record there. The asymmetric cell inserts the
sheet between an existing $H$ sample and an existing $E$ sample and adds no
nodes. The tight asymmetric cell adds one electric node on one side of the
cut and one magnetic node on the other, and it is the placement that paper
found to converge fastest. Milestone~5 codes~\eqref{eq:sc-h}--\eqref{eq:sc-amp-p}
and checks $R$ and $T$ against~\eqref{eq:R} and~\eqref{eq:T}. On the
validation grid the half-cell offset is the leading error, and refining
the cell drops it, so this placement stays. The tight asymmetric cell
remains the recorded alternative. The coded cell appends $H^\pm$ after
the electric unknowns. It does not paint the conductance~\eqref{eq:sigma-sheet}.

A constant real $\chi$, independent of $\omega$, is a convenient FDFD value
and a legal reduction of~\eqref{eq:sc-h}--\eqref{eq:sc-e}. It is not a causal
broadband model. Broadband FDTD uses a Lorentz oscillator for each resonance,
\begin{align}
\chi_{ee}(\omega) &= \frac{\omega_{ep}^2}{\omega_{e0}^2-\omega^2+j\alpha_e\omega},
&
\partial_t^2 P_z + \alpha_e\partial_t P_z + \omega_{e0}^2 P_z
&= \varepsilon_0\omega_{ep}^2 E_{z,\mathrm{av}},
\label{eq:lorentz-e}\\
\chi_{mm}(\omega) &= \frac{\omega_{mp}^2}{\omega_{m0}^2-\omega^2+j\alpha_m\omega},
&
\partial_t^2 M_x + \alpha_m\partial_t M_x + \omega_{m0}^2 M_x
&= \omega_{mp}^2 H_{x,\mathrm{av}}.
\label{eq:lorentz-m}
\end{align}
Several poles superpose. The auxiliary variables are marched with the local
cell solve. This is the dispersive extension, after the constant-$\chi$ sheet
matches~\eqref{eq:R} and~\eqref{eq:T}.

\subsection{What waits until the uniform sheet matches}
\label{sec:gstc-later}

Spatially varying $\chi(\mathbf{r})$ is a different $(4\times 4)$ on each cut
face, with the same jumps. A Bloch cell multiplies the periodic wrap by
$e^{-j\mathbf{k}_B\cdot\mathbf{L}}$ and reports $R$ and $T$ in the Floquet
orders. Spatially dispersive $\chi(k_t)$ is a separate operator, checked
against MaxwellMOM's multi-$k_t$ reductions, and it is not a local $4\times 4$.
None of these three is in this slice.

\section{Software architecture}
\label{sec:architecture}

\subsection{Modules}
\label{sec:modules}

The modules are the layout the tests exercise.

\begin{center}
\begin{tabular}{@{}ll@{}}
\toprule
Module & Responsibility\\
\midrule
\texttt{grid.yee2d} & Counts, spacings, polarization, boundary, shapes, coordinates\\
\texttt{materials.volume} & Uniform, staircase, and mapped $\varepsilon,\mu,\sigma$\\
\texttt{materials.dispersion} & Debye and Lorentz laws, auxiliaries $J_p$\\
\texttt{materials.conductors} & Staircase PEC and PMC masks\\
\texttt{operators.curl2d} & \texttt{curl\_e}, \texttt{curl\_h}, array curls, DOF order, optional $1/s_w$\\
\texttt{operators.pml} & Stretch~\eqref{eq:stretch}--\eqref{eq:sigma-max}; 2D auxiliary and 3D FDFD\\
\texttt{drivers.fdtd2d} & Leapfrog step, $J$ and $J_p$, invariant~\eqref{eq:invariant}\\
\texttt{drivers.fdfd2d} & System~\eqref{eq:fdfd}, Dirichlet lift, optional sheet\\
\texttt{analytics.cavity2d} & Frequency~\eqref{eq:semi} and mode shapes~\eqref{eq:mode-tm}, \eqref{eq:mode-te}\\
\texttt{analytics.dispersion} & Frequency~\eqref{eq:discrete} and the three-sample recurrence\\
\texttt{analytics.fresnel} & Slab fields~\eqref{eq:fresnel-field}, \eqref{eq:fresnel-te}\\
\texttt{analytics.mie2d} & Cylinder series~\eqref{eq:mie-tm}, \eqref{eq:mie-te}\\
\texttt{analytics.green2d} & Outgoing line current, $H_0^{(2)}$\\
\texttt{grid.yee3d} & 3D counts, spacings, boundary, shapes, coordinates\\
\texttt{materials.volume3d} & Uniform law and a staircase ball\\
\texttt{operators.curl3d} & 3D curls and array curls, optional $1/s_w$\\
\texttt{drivers.fdfd3d} & System~\eqref{eq:fdfd}, optional PML\\
\texttt{analytics.cavity3d} & Frequency~\eqref{eq:semi} plus the $z$ term\\
\texttt{analytics.mie3d} & Sphere series with $h_n^{(2)}$\\
\texttt{materials.sheets} & Resistive $\sigma_{\mathrm{eff}}$, tangential and normal $\chi$\\
\texttt{operators.gstc} & Symmetric cell, append $H^\pm$\\
\texttt{analytics.sheets} & $R,T$ from~\eqref{eq:R}, \eqref{eq:T}\\
\bottomrule
\end{tabular}
\end{center}

\texttt{YeeGrid2D} owns $n_x$, $n_y$, $\Delta x$, $\Delta y$, the polarization,
and the boundary kind. Both drivers read shapes from it. A validation case is
run in time or in frequency by constructing the other driver on the same
object. \texttt{YeeGrid3D} owns the three counts, the three spacings, and the
same boundary kind. The 3D grid is its own object: every component that the
stencil staggers in $z$ carries that stagger. Its FDFD driver reads the same
elimination, \eqref{eq:fdfd}, on $(E_x,E_y,E_z)$. An optional PML rebuilds
those curls with the same $1/s_w$ as~\eqref{eq:stretch}. A staircase ball
assigns one isotropic permittivity to every sample inside the closed sphere.

\texttt{UniformIsotropic.sample\_tm} and \texttt{sample\_te} take the grid and
optional \texttt{omega} and \texttt{state} arguments. Milestone~1 passes
neither. A dispersive or externally modulated law implements the same methods
and returns arrays of the same shapes. The stepper does not learn about
temperature.

A resistive sheet is not a new unknown. \texttt{ResistiveSheet} adds
$1/(Z_s\Delta y)$ to $\sigma$ on the tangential electric samples of the
Ampere cell in Sec.~\ref{sec:impedance}. Both drivers then use that
$\sigma$ unchanged. The symmetric cell of Sec.~\ref{sec:symmetric-cell}
replaces the magnetic sample on the cut with the pair $H^-$, $H^+$.
\texttt{operators.gstc} appends that pair, and the solve returns the
electric prefix. PEC remains a boundary kind rather
than a material, because it removes unknowns instead of filling
coefficients. An embedded PEC or PMC applies that removal to interior
samples, through \texttt{materials.conductors}.

\subsection{Relation to MaxwellMOM}
\label{sec:mom}

MaxwellMOM is a surface integral code: RWG bases, Green's functions, and dense
blocks. MaxwellFD does not import it. The shared pieces are the convention and
the analytic targets.

\begin{itemize}
\item $e^{j\omega t}$, and the four SI constants, copied numerically.
\item Sheet reflection and transmission,~\eqref{eq:R} and~\eqref{eq:T},
including the oblique factors and the $\chi_{mm}^{nn}$ correction, which are
the formulas behind \texttt{rt\_from\_chi}.
\item The map~\eqref{eq:chi-from-z} from a surface impedance to $\chi_{ee}$.
\item Mie and Fresnel results as external checks, not as shared functions.
\end{itemize}

A number computed by MaxwellMOM for a configuration both codes can run (a flat
sheet, a periodic cell) is a regression oracle. It is stored as a literal in a
test, not read from the other repository at runtime.

\subsection{Milestone map}
\label{sec:milestones}

\begin{enumerate}
\item \textbf{This slice.} 2D grid, both drivers, PEC cavities, periodic
dispersion, invariant~\eqref{eq:invariant}.
\item Volumetric $\varepsilon(\mathbf{r})$, then $\varepsilon(\omega)$: complex
permittivity in FDFD, Debye or Lorentz auxiliaries in FDTD. Slab versus
Fresnel, cylinder versus 2D Mie. External maps enter only as arrays the law
samples.
\item PML from Sec.~\ref{sec:pml}, embedded PEC and PMC objects, an open
PEC cylinder, and an open dielectric cylinder, both against the Mie
series. The 3D cell, a rectangular PEC cavity, and an open dielectric
sphere small enough to factor are in the code.
\item Resistive sheet, Sec.~\ref{sec:impedance}. The conductance
\eqref{eq:sigma-sheet} is in the code. It is compared with~\eqref{eq:R},
\eqref{eq:T}, and one stored MaxwellMOM value.
\item Symmetric GSTC cell, Sec.~\ref{sec:symmetric-cell}. Tangential
$\chi_{ee}$ and $\chi_{mm}$, and the normal susceptibility
$\chi_{mm}^{nn}$, are in the code and are compared
with~\eqref{eq:R} and~\eqref{eq:T}. The bianisotropic terms
in~\eqref{eq:bianiso} remain later.
\item Dipole, Gaussian beam, near-to-far, Bloch phase on the periodic wrap.
\item Profiling, Numba, iterative FDFD, GPU. Spatially varying and spatially
dispersive $\chi$ remain later.
\end{enumerate}

Acceleration does not change the curls. Numba and CuPy, when they arrive,
replace the array update and the sparse matrix-vector product. The direct
factorization remains the 2D reference.

\begin{thebibliography}{9}
\bibitem{yee1966}
K.~S. Yee, ``Numerical solution of initial boundary value problems involving
Maxwell's equations in isotropic media,'' \emph{IEEE Trans.\ Antennas
Propag.}, 1966.
\bibitem{gedney2000}
J.~A. Roden and S.~D. Gedney, ``Convolution PML (CPML): An efficient FDTD
implementation of the CFS--PML for arbitrary media,'' \emph{Microwave Opt.\
Technol.\ Lett.}, 2000.
\bibitem{smy2017}
T.~J. Smy, S. Stewart, and S. Gupta, ``Integrated generalized sheet transition
conditions in a Yee-cell based FDTD simulation of electromagnetic
metasurfaces,'' arXiv:1706.10136, 2017.
\bibitem{vahabzadeh2017}
Y. Vahabzadeh, N. Chamanara, and C. Caloz, ``Generalized sheet transition
condition FDTD simulation of metasurface,'' arXiv:1701.08760, 2017.
\end{thebibliography}

\end{document}
